\documentclass[11pt,a4paper]{article}
\usepackage[utf8]{inputenc}
\usepackage{amsmath, amssymb, amsfonts}
\usepackage{booktabs}
\usepackage{geometry}
\geometry{a4paper, margin=0.85in}
\usepackage{hyperref}
\usepackage{xcolor}

\hypersetup{
    colorlinks=true,
    linkcolor=blue!70!black,
    citecolor=blue!70!black,
    urlcolor=blue!70!black
}

\title{\textbf{Mathematical Research Evidence: Three Methodological Case Studies from BrainKAN} \\
\large Dossier on Mathematical Modeling, Hypothesis Testing, and Inferential Boundaries}
\author{\textbf{Xin Qi} \\ Department of Computing / Neuroimaging AI Benchmark Release}
\date{October 2026}

\begin{document}
\maketitle

\begin{abstract}
This dossier documents three concrete methodological case studies from the BrainKAN benchmark, formulated to demonstrate core research competencies in: (1) mathematical modeling and function space analysis (operational metric definition and synthetic counterexample construction); (2) probability, hypothesis testing, and inferential boundaries (exchangeability in paired permutation null models); and (3) sensitivity analysis and statistical critique (estimand duality and directional sign-inversion under domain support constraints). Rather than highlighting high classification accuracy ($93\%$) or standard B-spline recursion formulas, these cases demonstrate the ability to interrogate operational mathematical definitions, diagnose inferential failure modes, and demarcate the exact boundaries of empirical scientific claims.
\end{abstract}

\vspace{0.5em}
\hrule
\vspace{1em}

\section{Case 1: Mathematical Modeling \& Metric Definition: Reference Domain Sensitivity of DFL}

\subsection{Mathematical Definition}
Let $f: \mathbb{R} \to \mathbb{R}$ denote a univariate transformation (e.g., an edge-specific B-spline function in a Kolmogorov--Arnold Network), evaluated on a discrete probe grid $D = \{x_1, \dots, x_K\} \subset \mathbb{R}$ with $K \ge 3$ distinct points. To quantify the deviation of $f$ from the subspace of affine linear polynomials $\mathcal{P}_1 = \{ax + b : a, b \in \mathbb{R}\}$, we compute the best linear approximation in the $\ell_2(D)$ norm via ordinary least squares:
\begin{equation}
\hat{y}(x) = a^* x + b^*, \quad (a^*, b^*) = \arg\min_{a, b \in \mathbb{R}} \sum_{i=1}^K [f(x_i) - (a x_i + b)]^2
\end{equation}
The \textbf{Deviation-From-Linearity ($\operatorname{DFL}$)} is defined as the normalized residual sum of squares:
\begin{equation}
\boxed{
\operatorname{DFL}(f; D) = \frac{\min_{a, b} \sum_{i=1}^K [f(x_i) - (a x_i + b)]^2}{\sum_{i=1}^K [f(x_i) - \bar{f}_D]^2} = \frac{\mathrm{SS}_{\mathrm{res}}}{\mathrm{SS}_{\mathrm{tot}}} = 1 - R^2
}
\end{equation}
where $\bar{f}_D = \frac{1}{K} \sum_{i=1}^K f(x_i)$.

\subsection{Analytical Nuances \& Non-Degeneracy Conditions}
\begin{enumerate}
    \item \textbf{Output Scale and Shift Invariance:} For any affine transformation $g(x) = \alpha f(x) + \beta$ ($\alpha \ne 0, \beta \in \mathbb{R}$), $\operatorname{DFL}(\alpha f + \beta; D) = \operatorname{DFL}(f; D)$. The metric is dimensionless and scale-invariant with respect to output magnitude.
    \item \textbf{Directionality:} $\operatorname{DFL}$ measures the fraction of total variance unexplained by the affine model. Higher values indicate greater departure from linearity ($\operatorname{DFL} = 0$ corresponds to perfect linearity on $D$; $\operatorname{DFL} = 1$ indicates zero linear correlation).
    \item \textbf{Degeneracy Convention:} The metric assumes $\mathrm{SS}_{\mathrm{tot}} > 0$. If $f$ is constant on $D$ ($f(x_i) \equiv c$), $\mathrm{SS}_{\mathrm{tot}} = 0$, yielding an indeterminate $0/0$. In our implementation, constant functions are assigned $\operatorname{DFL}(f; D) = 0$ by operational convention.
    \item \textbf{Discrete Grid Limitation:} $\operatorname{DFL}(f; D) = 0$ guarantees collinearity only at the discrete evaluation points $\{x_i\}_{i=1}^K$; it does not mathematically guarantee that the continuous spline is affine between probe points.
\end{enumerate}

\subsection{Counterexample \& Synthetic Domain Sensitivity}
To evaluate whether $\operatorname{DFL}$ behaves as an intrinsic property of $f$, we audited it using a controlled synthetic experiment (Exp 2):
\begin{itemize}
    \item \textbf{Setup:} A static ground-truth function $f(x) = \tanh(x)$ was evaluated across two reference domains: $D_1 = [-2, 2]$ and $D_2 = [0, 2]$.
    \item \textbf{Finding:} Evaluating the \textit{exact same function} on $D_1$ versus $D_2$ produced an apparent metric jump of $\Delta \operatorname{DFL} \approx 0.0816 \pm 0.0084$ ($>40\times$ larger than typical task differences in fMRI networks).
    \item \textbf{Crucial Technical Qualification:} $D_1 = [-2, 2]$ and $D_2 = [0, 2]$ differ in both center ($0$ vs. $1$) and width/support length ($4$ vs. $2$). Therefore, this discrepancy cannot be characterized as a pure translational effect; it reflects the joint sensitivity of the metric to domain location and scale.
\end{itemize}
\textbf{Competence Demonstrated:} Function space understanding, projection operator geometry, and constructing synthetic counterexamples to audit operational metrics.

\vspace{1em}
\hrule
\vspace{1em}

\section{Case 2: Probability, Hypothesis Testing \& Scientific Inference: Permutation Null Design and Boundaries}

\subsection{Formal Null Formulation \& Exchangeability}
In task-fMRI connectomics, each human subject underwent both a 0-Back working memory block and a 2-Back working memory block in a paired design. Let $\mathcal{D} = \{(G_{i,0}, G_{i,1})\}_{i=1}^N$ denote the paired graph datasets for $N = 100$ subjects.
\begin{itemize}
    \item \textbf{Null Hypothesis ($H_0$):} Task condition labels $Y \in \{0, 1\}$ are exchangeable within each subject $i$:
    \begin{equation}
    \mathcal{L}\left((G_{i,0}, G_{i,1}) \mid Y\right) = \mathcal{L}\left((G_{i,0}, G_{i,1}) \mid \pi_i(Y)\right), \quad \forall i \in \{1, \dots, N\}
    \end{equation}
    \item \textbf{Invariance Structure:} Within-subject sign-flips $\pi_i \in \{(0, 1), (1, 0)\}$ with equal probability $0.5$ preserve individual anatomy ($48$ ROIs) and functional connectome topology ($240$ edges) under the null, breaking strictly the association between learned weights and task supervision.
\end{itemize}

\subsection{Test Statistic \& Exact Empirical p-value}
With test statistic $T = \frac{1}{|E|} \sum_{e \in E} \operatorname{DFL}_e$ ($T_{\mathrm{obs}} = 0.17895$ for the real model), evaluating $B = 500$ independent paired permutation null models $\{T_b\}_{b=1}^{500}$ (Exp 1) yields the exact non-parametric permutation $p$-value:
\begin{equation}
\boxed{
p = \frac{1 + \#\{T_b \ge T_{\mathrm{obs}}\}}{B + 1} = \frac{1 + 2}{500 + 1} = \frac{3}{501} \approx 0.005988 \approx 0.00599
}
\end{equation}
The 500 null values yielded empirical mean $\bar{T}_{\mathrm{null}} = 0.17733$ and sample standard deviation $\sigma_{\mathrm{null}} = 0.000585$, establishing a Monte Carlo standard error of $\mathrm{SE}_{\mathrm{MC}} = \frac{0.000585}{\sqrt{500}} \approx 0.000026$.

\subsection{Scientific Demarcation \& Inferential Boundaries}
Rather than stopping at reporting statistical significance ($p \approx 0.006$), we explicitly articulate the inferential boundaries:
\begin{enumerate}
    \item The $p$-value quantifies the probability of observing an effect of this magnitude under the specific null hypothesis of exchangeable task labels, conditional on the fixed cohort sample and training regime.
    \item Rejecting $H_0$ demonstrates that the learned network weights encode information non-exchangeable with respect to task labels. However, because overparameterized neural networks are non-identifiable in parameter space, \textbf{this statistical rejection does not prove that individual learned edge splines correspond to in-vivo biological mechanisms}.
    \item The single-seed pooled-domain permutation test in Exp 1 evaluates a specific model instance on a fixed dataset; it does not replace out-of-distribution validation on external clinical cohorts.
\end{enumerate}
\textbf{Competence Demonstrated:} Exchangeability in paired experimental designs, finite-sample permutation testing, Monte Carlo error quantification, and intellectual restraint in scientific inference.

\vspace{1em}
\hrule
\vspace{1em}

\section{Case 3: Sensitivity Analysis \& Estimand Duality: The Sign Inversion under Common Support}

\subsection{The Critical Phenomenon: Directional Inversion of Results}
The most compelling methodological finding is that the empirical conclusion \textbf{inverts in sign} depending on whether the evaluation domain accounts for extrapolation across condition boundaries.

In Exp 5, across $S = 10$ independent random optimization seeds ($w_0 \sim \mathcal{S}$), we evaluated the paired excess curvature $\Delta_s = T(Y; s) - \frac{1}{15} \sum_{b=1}^{15} T(\pi_b(Y); s)$:

\begin{table}[h]
\centering
\small
\begin{tabular}{llcccc}
\toprule
\textbf{Evaluation Estimand} & \textbf{Mathematical Domain} & \textbf{Mean Diff ($\bar{\Delta}$)} & \textbf{Seed Consistency} & \textbf{Student's $t$ ($df=9$)} & \textbf{One-Sided $p$} \\
\midrule
\textbf{Pooled Domain ($\mathcal{S}_{\mathrm{pooled}}$)} & $\text{Quantile}_{1\%-99\%}([X_0, X_1])$ & $\mathbf{-0.000583}$ ($\pm 0.00100$) & \textbf{1 / 10 Positive} (Seed 42) & $t = -1.849$ & $p = 0.9513$ \\
\textbf{Strict Common Support ($\mathcal{S}_{\mathrm{common}}$)} & $\operatorname{supp}(X_0) \cap \operatorname{supp}(X_1)$ & $\mathbf{+0.000896}$ ($\pm 0.00055$) & \textbf{9 / 10 Positive} & $t = +5.186$ & $p = 0.0003$ \\
\bottomrule
\end{tabular}
\end{table}

\subsection{Mathematical \& Methodological Diagnosis}
\begin{enumerate}
    \item \textbf{Extrapolation on the Pooled Grid:} Across 48 ROIs on HCP fMRI data, the Jaccard overlap between 0-Back and 2-Back empirical distributions averages $38.52\% \pm 5.08\%$ (covering $43.47\%$ of the pooled span). On $\mathcal{S}_{\mathrm{pooled}}$, splines are evaluated outside their empirical training range; boundary polynomial oscillation dominates subtle task signals, producing negative expectations for 9 out of 10 initializations.
    \item \textbf{Isolation of Mutual Support:} Restricting evaluation to the strict intersection $\mathcal{S}_{\mathrm{common}} = [\max(q_{1,0}, q_{1,1}), \min(q_{99,0}, q_{99,1})]$ filters out extrapolation noise. Under this invariant domain, 9 out of 10 seeds consistently exhibit positive excess curvature.
\end{enumerate}

\subsection{Critical Thinking \& Inferential Boundaries}
This case exemplifies research maturity: an apparently positive result ($p = 0.006$ on Seed 42) collapsed across seeds on the pooled domain, but emerged consistently when the domain was bounded to common support.
\begin{itemize}
    \item \textbf{Optimization Stochasticity vs. Population Generalization:} The paired $t$-statistic across seeds ($t = 5.186, p = 0.0003$) evaluates \textit{initialization consistency conditional on the fixed cohort sample and fixed 15 permutation vectors}. It does not claim generalizability to unseen biological populations ($\mathcal{D} \sim \mathcal{P}$).
    \item \textbf{Shared Permutation Vectors \& Baseline Error:} The 10 seeds shared the same 15 permutation vectors, leaving common baseline error unmodeled across seeds ($\mathrm{SE}_{\mathrm{MC}} \approx 0.00015$ by illustrative order-of-magnitude proxy).
    \item \textbf{Label-Dependent Domain Selection:} The common support interval was constructed using empirical condition distributions. While holding this grid fixed provides measurement comparability, formal finite-sample type I error rate control under label-dependent domain selection remains an open statistical question.
\end{itemize}
\textbf{Competence Demonstrated:} Diagnosing sign inversions driven by support boundaries, estimand demarcation, and refusing to overclaim optimization stability as population-level proof.

\vspace{1em}
\hrule
\vspace{1em}

\section*{Summary Matrix of Methodological Competencies}
\begin{table}[h!]
\centering
\small
\begin{tabular}{lll}
\toprule
\textbf{Competency Area} & \textbf{Methodological Challenge \& Solution} & \textbf{Demonstrated Skill} \\
\midrule
\textbf{Mathematical Modeling} & Formalized $\operatorname{DFL} = \frac{\mathrm{SS}_{\mathrm{res}}}{\mathrm{SS}_{\mathrm{tot}}} = 1-R^2$; audited on shifted domains & Function space analysis, projection operators, counterexamples \\
\textbf{Probability \& Statistics} & Paired permutation null via within-subject sign flips ($B=500, p \approx 0.00599$) & Exchangeability, non-parametric inference, Monte Carlo bounds \\
\textbf{Critical Thinking} & Discovered sign flip ($\bar{\Delta}: -0.00058 \to +0.00090$) between pooled/common grids & Estimand definition, domain sensitivity, inferential discipline \\
\bottomrule
\end{tabular}
\end{table}

\end{document}
